\documentclass[10pt,a4paper]{article}
\usepackage[a4paper,left=0.78in,right=0.78in,top=0.72in,bottom=0.72in]{geometry}
\usepackage{fontspec}
\setmainfont{Noto Serif}
\usepackage{xcolor}
\usepackage{colortbl}
\definecolor{navy}{RGB}{21,48,75}
\usepackage{tikz}
\usepackage{fancyhdr}
\usepackage{array}
\usepackage{tabularx}
\usepackage{enumitem}
\usepackage{ragged2e}
\usepackage{hyperref}
\hypersetup{colorlinks=true,linkcolor=navy,urlcolor=navy}
\setlength{\parindent}{0pt}
\setlength{\parskip}{2.5pt}
\setlength{\headheight}{14pt}
\setlist[itemize]{leftmargin=1.5em,itemsep=1pt,topsep=1pt}
\newcommand{\pageframe}{%
  \begin{tikzpicture}[remember picture,overlay]
    \draw[navy,line width=0.8pt] ([xshift=0.23in,yshift=-0.23in]current page.north west) rectangle ([xshift=-0.23in,yshift=0.23in]current page.south east);
  \end{tikzpicture}%
}
\fancypagestyle{main}{%
  \fancyhf{}
  \fancyhead[L]{\textcolor{navy}{\small 1008 -- Foundational IT Skills \;|\; Semester 1}}
  \fancyfoot[L]{\textcolor{gray}{\footnotesize STATE BOARD OF TECHNICAL EDUCATION KERALA \;|\; DIPLOMA CURRICULUM REVISION 2026}}
  \fancyfoot[R]{\textcolor{gray}{\footnotesize \thepage}}
  \renewcommand{\headrulewidth}{0pt}
  \renewcommand{\footrulewidth}{0pt}
}
\AddToHook{shipout/background}{\pageframe}
\newcommand{\expheading}[2]{%
  \vspace{0.1cm}\noindent\textbf{\large EXPERIMENT NO: #1}\hfill\textbf{Date: #2}\par
}
\newcommand{\exptitle}[1]{%
  \begin{center}\vspace{-2mm}\textcolor{navy}{\fontsize{20}{22}\selectfont\bfseries #1}\end{center}
  \vspace{1mm}\hrule height 0.5pt\vspace{3mm}
}
\begin{document}
\pagestyle{main}
\thispagestyle{empty}
\begin{center}
\vspace*{1.0cm}
{\textcolor{navy}{\fontsize{18}{21}\selectfont\bfseries STATE BOARD OF TECHNICAL EDUCATION}}\\[-2mm]
{\textcolor{navy}{\fontsize{18}{21}\selectfont\bfseries KERALA}}\\[1mm]
{\textcolor{navy}{\large\bfseries DIRECTORATE OF TECHNICAL EDUCATION}}\\
{\large GOVERNMENT OF KERALA}

\vspace{1.3cm}
{\textcolor{navy}{\fontsize{23}{27}\selectfont\bfseries Diploma in Engineering}}\\[5mm]
\colorbox{navy}{\parbox[c][0.48in][c]{3.5in}{\centering\color{white}\Large\bfseries SEMESTER 1}}

\vspace{1.2cm}
{\textcolor{navy}{\fontsize{22}{26}\selectfont\bfseries LABORATORY MANUAL}}

\vspace{1.0cm}
\noindent\begin{tabular}{|>{\centering\arraybackslash\bfseries}p{1.15in}|>{\raggedright\arraybackslash\color{white}\bfseries}p{4.55in}|}\hline
\textcolor{navy}{\Large 1008} & \cellcolor{navy}\Large FOUNDATIONAL IT SKILLS\\\hline
\end{tabular}

\vspace{1.0cm}
{\textcolor{navy}{\large\itshape For Diploma Curriculum Revision 2026}}

\vspace{1.0cm}
{\textcolor{navy}{\Large\bfseries STUDENT COPY}}
\end{center}
\clearpage
\thispagestyle{main}
\begin{center}{\textcolor{navy}{\fontsize{20}{24}\selectfont\bfseries INDEX}}\end{center}
\vspace{5mm}
\renewcommand{\arraystretch}{1.35}
\begin{tabularx}{\textwidth}{|c|X|c|c|c|}\hline
\textbf{Exp. No.} & \textbf{Name of Experiment} & \textbf{Date} & \textbf{Page No.} & \textbf{Faculty Signature}\\\hline
1 & Hardware \& Software Exploration & & & \\\hline
2 & File And Folder Management In Linux Gui & & & \\\hline
3 & Web \& Connectivity & & & \\\hline
4 & Professional Documentation & & & \\\hline
5 & Spreadsheet Basics & & & \\\hline
6 & Presentation Design & & & \\\hline
7 & Cloud Collaboration & & & \\\hline
8 & Data Analysis & & & \\\hline
9 & Data Visualization & & & \\\hline
10 & Ai Slide Generation & & & \\\hline
11 & Generative Pre-Trained Transformer (Gpt) Guided Learning & & & \\\hline
12 & Graphic Design & & & \\\hline
\end{tabularx}
\clearpage
\clearpage
\expheading{1}{\underline{\hspace{1.8cm}}}
\exptitle{HARDWARE \& SOFTWARE EXPLORATION}
\medskip\noindent\textbf{PRACTICAL SIGNIFICANCE}\par
\noindent To identify the major hardware components of a computer system and familiarize students with\par
\noindent system software and application software available on a personal computer.\par

\medskip\noindent\textbf{RELEVANT PROGRAM OUTCOMES(POS)}\par
\noindent PO1: Apply knowledge of mathematics, natural sciences, engineering fundamentals, and an\par
\noindent engineering specialization (as specified in DK1 to DK4) to a wide range of practical procedures and\par
\noindent practices.\par

\medskip\noindent\textbf{RELEVANT COURSE OUTCOMES(COS)}\par
\noindent CO1 : Identify the basic functions and features of Computers, Operating Systems, and Internet\par
\noindent applications.\par
\noindent THEORY Write the minimum theory required for performing the experiment.\par

\medskip\noindent\textbf{Hardware}\par
\noindent Hardware refers to the physical components of a computer system that can be seen and touched,\par

\noindent Major hardware components include:\par
\noindent 1. CPU (Central Processing Unit)\par
\noindent\textbullet\; CPU is called the "brain of the computer."\par
\noindent\textbullet\; It processes instructions and performs calculations.\par
\noindent\textbullet\; It controls all operations of the computer system.\par
\noindent\textbullet\; Main components: Arithmetic Logic Unit (AIU), Control Unit (CU), Registers Functions:\par
\noindent\textbullet\; Executes instructions\par
\noindent\textbullet\; Processes data\par
\noindent\textbullet\; Controls system operations\par
\noindent AIN\par
\medskip\noindent\textbf{SCHEME}\par
\noindent 2. RAM (Random Access Memory)\par
\noindent\textbullet\; RAM is temporary memory used by the computer while executing programs.\par
\noindent\textbullet\; Data stored in RAM is lost when power is switched off.\par

\noindent Functions:\par
\noindent\textbullet\; Stores currently running programs\par
\noindent\textbullet\; Stores temporary data\par

\noindent Characteristics:\par
\noindent\textbullet\; Volatile memory\par
\noindent\textbullet\; Faster than secondary storage\par
\noindent 3. HDD (Hard Disk Drive)\par
\noindent\textbullet\; HDD is a secondary storage device.\par
\noindent\textbullet\; Stores operating systems, files, software, and user data permanently. Characteristics:\par
\noindent\textbullet\; Non-volatile storage\par
\noindent\textbullet\; Mechanical storage device\par
\clearpage
\noindent 4. SSD (Solid State Drive)\par
\noindent\textbullet\; SSD is a storage device using flash memory.\par
\noindent\textbullet\; It has no moving parts.\par

\noindent Advantages over HDD:\par
\noindent\textbullet\; Faster access speed\par
\noindent\textbullet\; Less power consumption\par
\noindent\textbullet\; More durable\par
\noindent 5. Peripheral Ports\par
\noindent Peripheral ports are interfaces used to connect external devices to the computer.\par

\noindent ‘Common ports:\par
\noindent Port Purpose USB Port Connect keyboard, mouse, pen drive, printer\par
\noindent HDMI Port Connect monitor/projector\par
\noindent VGA Port Connect display devices Ethernet Port ‘Network connection ‘Audio Port Connect\par
\noindent speakers/headphones Power Port ‘Supply power\par
\noindent Type-C Port Data transfer and charging\par

\medskip\noindent\textbf{Software}\par
\noindent Software is a collection of programs and instructions that directs the computer to perform specific\par
\noindent tasks.\par

\noindent Software is classified into:\par
\noindent 1. System Software\par
\noindent System software controls the operation of the computer and manages hardware resources. Examples:\par
\noindent\textbullet\; Operating System\par
\noindent\textbullet\; Device Drivers\par
\noindent\textbullet\; Language Translators\par
\noindent\textbullet\; Utility Software Examples available on PC:\par
\noindent\textbullet\; Windows\par
\noindent\textbullet\; Linux\par
\noindent\textbullet\; Antivirus software\par
\noindent\textbullet\; Disk Cleanup utility\par
\noindent\textbullet\; Device drivers Functions:\par
\noindent\textbullet\; Manages hardware\par
\noindent\textbullet\; Provides user interface\par
\noindent\textbullet\; Controls system resources\par
\noindent 2. Application Software\par
\noindent Application software performs specific user tasks.\par

\noindent Examples:\par
\noindent Software Purpose MS Word ‘Document creation MS Excel ‘Spreadsheet operations ‘MS PowerPoint\par
\noindent Presentation preparation\par
\noindent Google Chrome Web browsing\par

\noindent VLC Media Player                                Playing media files\par

\noindent Paint Image editing and drawing,\par
\noindent Calculator Calculations\par
\clearpage
\noindent Sl.No.                      \_ Name of resource         Computer System               1 per student\par
\noindent Suggested specification    Desktop/Laptop with\par
\noindent Qua 1                      Ubuntu\par

\noindent Linux\par

\noindent 2                                                      Monitor 15 inch or above 1\par

\noindent 3                                                      Keyboard Standard USB keyboard 1\par

\noindent 4                                                      Mouse Standard USB mouse 1\par

\noindent 5                                                      RAM Module Any available module for 1\par
\noindent demonstration\par

\noindent 6                                                      HDD/ssD Any available drive for 1 demonstration\par

\noindent 7                                                      Internet Connection Broadband / Wi-Fi 1 per lab\par

\noindent 8—                                     Ubuntu Linux OS Version 22.04     Installed\par
\noindent LTS or above\par

\noindent 9                                                      Web Browser Mozilla Firefox / Google Installed\par

\noindent Chrome\par
\medskip\noindent\textbf{(TEACHERS AND ERY ~ DIPLOMA 2}\par
\noindent ME\par

\medskip\noindent\textbf{PROCEDURE}\par
\noindent Hardware Identification\par
\noindent 1. Observe the computer system carefully.\par
\noindent 2. Identify CPU.\par
\noindent 3. Identify RAM.\par
\noindent 4. Identify HDD/SSD.\par
\noindent 5. Observe peripheral ports available on the computer.\par
\noindent 6. Note down the purpose of each port. @ Software Identification\par
\noindent 1. Switch on the computer.\par
\noindent 2. Open installed applications.\par
\noindent 3. Identify system software.\par
\noindent 4, Identify application software.\par
\noindent 5. Record observations.\par

\medskip\noindent\textbf{RESULTS}\par
\noindent Thus the major hardware components and software components of a personal computer were\par
\noindent identified and studied successfully.\par
\medskip\noindent\textbf{PRACTICAL RELATED QUESTPRACTICAL RELATED QUESTPRACTICAL RELATED QUESTIONS}\par
\noindent 1. What is CPU?\par
\noindent 2. Expand RAM.\par
\noindent 3. What is the difference between RAM and ROM?\par
\noindent 4. What is HDD?\par
\noindent 5. What is SSD?\par
\noindent 6. Name any two peripheral ports.\par
\clearpage
\noindent 7. What is system software?\par
\noindent 8. What is application software?\par
\noindent 9. Give examples of system software.\par
\noindent 10. Give examples of application software.\par

\medskip\noindent\textbf{SELF-LEARNING COMPONENTS}\par
\noindent 1. Draw or Take a picture of the ports of your Lab PC and label the ports of your PC.\par
\noindent 2. Compare the RAM, Processor, and Storage type (SSD vs. HDD) of your Tablet PC or any other\par
\noindent smart devices versus your lab PC. Prepare a comparison chart of the configurations.\par
\noindent 3. List any 5 system software and application software in your lab PC.\par
\noindent REFERENCES SI.No. Text Book/ Website ‘Author Publication\par

\noindent 1                           https:/www.geeksforgeek   GeeksforGeeks      GeeksforGeeks g/comput\par
\noindent s.or                                         er-science-fundame ntals/\par
\noindent computer-hardware/\par
\clearpage
\expheading{2}{\underline{\hspace{1.8cm}}}
\exptitle{FILE AND FOLDER MANAGEMENT IN}
\medskip\noindent\textbf{LINUX GUI}\par
\noindent PRACTICAL SIGNIFICANCE To create and manage files and folders in Linux using the Graphical\par
\noindent User Interface (GUI. PO1: Apply knowledge of mathematics, natural sciences, engineering\par
\noindent fundamentals, and an\par
\noindent engineering specialization (as specified in DK1 to DK4) to a wide range of practical procedures and\par
\noindent practices.\par

\medskip\noindent\textbf{RELEVANT COURSE OUTCOMES(COS)}\par
\noindent CO1 : Identify the basic functions and features of Computers, Operating Systems, and Internet\par
\noindent applications.\par
\noindent THEORY Write the minimum theory required for performing the experiment.\par

\noindent Linux Operating System\par
\noindent Linux is an open-source operating system that manages computer hardware and software resources.\par
\noindent It supports both command-line and graphical user interfaces.\par

\medskip\noindent\textbf{Features of Linux:}\par
\noindent Open-source operating system\par
\noindent Multi-user operating system\par
\noindent Multi-tasking capability\par
\noindent Secure and reliable\par
\noindent Supports GUI and command line interface\par

\noindent Graphical User Interface (GUI)\par
\noindent GUI allows users to interact with the computer using windows, icons, menus, and a mouse pointer.\par
\medskip\noindent\textbf{Advantages of GUI}\par
\noindent Easy to use\par
\noindent User-friendly\par
\noindent Does not require memorizing commands\par
\noindent Supports drag-and-drop operations\par

\noindent File\par
\noindent A file is a collection of information stored on a storage device.\par

\noindent Examples:\par
\noindent Text files, Images, Audio files, Video files, Documents\par

\noindent Folder (Directory)\par
\noindent A folder or directory is used to organize and store files and other folders.\par

\medskip\noindent\textbf{Functions of folders:}\par
\noindent files\par
\noindent Reduce confusion\par
\clearpage
\noindent Orga\par
\noindent Improve file management\par

\medskip\noindent\textbf{Common File Management Operations}\par
\noindent Create : Creating a new file or folder.\par
\noindent Rename : Changing the name of an existing file or folder. Copy: Creating a duplicate of a file or\par
\noindent folder.\par
\noindent Move: Changing the location of a file or folder.\par
\noindent Delete: Removing unwanted files or folders\par

\noindent RESOURCES REQUIRED          Name of resource          Computer System            1 per student\par
\noindent Sl.No.                      Suggested specification   Desktop/Laptop with\par
\noindent Quantity 1                Ubuntu\par

\noindent Linux 2                     Ubuntu Linux Os Version   \_ Installed 3              Files Application (File\par
\noindent 22.04 LTS or above                                   GNOME Files / Nautilus\par
\noindent Installed\par

\noindent Manager) (default in Ubuntu)\par

\noindent 4                                                     Text Editor Gedit / Mousepad Installed\par

\medskip\noindent\textbf{PROCEDURE}\par
\noindent Creating Folders\par
\noindent 1. Open File Manager.\par
\noindent 2. Navigate to Home folder.\par
\medskip\noindent\textbf{(TEACHERS AND ERY ~ DIPLOMA 2}\par
\noindent ME\par
\noindent 3. Right-click in empty space.\par
\noindent 4. Select New Folder.\par
\noindent 5. Create folder named ComputerLab.\par
\noindent Creating Files\par
\noindent 1. Open ComputerLab folder\par
\noindent 2. Create files:\par
\noindent © Student1.txt\par
\noindent © Student2.txt\par
\noindent © Student3.txt\par
\noindent Rename Operation\par

\noindent J. Rename Student                                     .txt as StudentData.txt.\par

\noindent Copy Operation\par
\noindent 1. Copy Student2.tst into another folder.\par
\noindent Move Operation\par
\noindent J. Move Student3.txt to Documents folder.\par
\noindent Delete Operation\par
\noindent 1. Delete unwanted files.\par
\clearpage
\medskip\noindent\textbf{LINUX COMMANDS RELATED TO FILE AND FOLDER MANAGEMENT ( Optional}\par
\medskip\noindent\textbf{Exercise For Computer allied branches)}\par
\noindent AIthough file and folder operations can be performed using the Graphical User Interface (GUI), Linux\par
\noindent also provides command-line utilities for performing the same tasks efficiently. Understanding these\par
\noindent commands helps students gain familiarity with Linux system administration and file management.\par
\noindent 1. pwd ~ Print Working Directory\par
\noindent Purpose: Displays the current directory path.\par
\noindent 2. Is — List Directory Contents\par
\noindent Purpose: Displays files and folders in the current directory. Usefiul Options: Is -I, ls -a\par
\noindent 3. mkdir — Create Directory\par
\noindent Purpose: Creates a new folder (directory).\par
\noindent 4. touch ~ Create Empty File\par
\noindent Purpose:Creates a new empty file.\par
\noindent 5. ed ~ Change Directory\par
\noindent Purpose:Changes the current working directory.\par
\noindent 6. cp Copy Files and Directories\par
\noindent Purpose:Creates a duplicate copy of a file or folder.\par
\noindent 7. mv ~ Move or Rename Files\par
\noindent Purpose:Moves files/folders or renames them.\par
\noindent 8. rm ~ Remove Files\par
\noindent Purpose:Deletes files.\par
\noindent 9, rmdir ~ Remove Empty Directory\par
\noindent Purpose: Deletes an empty folder.\par
\noindent 10. rm -r ~ Remove Directory and Contents\par
\noindent Purpose:Deletes a folder along with all files and subfolders\par
\noindent 11. cat — Display File Contents\par
\noindent Purpose: Displays the contents of a text file.\par
\noindent Displays the file contents on the terminal.\par
\noindent 12. clear ~ Clear Terminal Screen\par
\noindent Purpose:Clears all text displayed in the terminal window.\par

\noindent Practical Note:\par
\noindent The instructor of Computer allied branches may first demonstrate file and folder management using\par
\noindent the GUI File Manager and then show the equivalent Linux terminal commands so that students\par
\noindent understand both graphical and command-line approaches to file management. This helps bridge\par
\noindent GUI-based learning with Linux system administration concepts.\par

\medskip\noindent\textbf{RESULTS}\par
\noindent File and folder management operations such as creating, renaming, copying, moving, and deleting\par
\noindent files and folders were successfully performed using the Linux Graphical User\par
\noindent Interface (GUI). The student gained familiarity with the Linux desktop environment and file\par
\noindent management tools. A report containing the steps performed was prepared and successfully uploaded\par
\noindent through Google Classroom.\par
\clearpage
\medskip\noindent\textbf{PRACTICAL RELATED QUESTPRACTICAL RELATED QUESTPRACTICAL RELATED QUESTIONS}\par
\noindent What is Linux? Expand GUI. What is a file? What is a folder? What is File Manager? How do you\par
\noindent create a folder? What is the difference between copy and move? What is rename operation? What is\par
\noindent the purpose of Trash?\par
\noindent 10. Name two Linue desktop environments.\par

\medskip\noindent\textbf{SELF-LEARNING COMPONENTS}\par
\medskip\noindent\textbf{SP ENAWAWNS}\par
\noindent 1. Use the GUI to create a directory tree for your semester subjects, using subject codes as ‘folder\par
\noindent names and creating separate sub-folders to organize notes for each module.\par
\noindent 2. Create folders and files for the given scenario\par
\noindent You are preparing for a weekend trip. Your "home base" is your main user folder.Create a new folder\par
\noindent called Backpack.Go inside your Backpack folder and create two folders: Food and Tools.Inside the\par
\noindent Tools folder, create a file called flashlight.txt and a file called map.txt. You realize the map.txt\par
\noindent shouldn't be in Tools. Move it into the Food folder instead. You want a second flashlight just in case.\par
\noindent Copy the flashlight.txt file and name the copy spare\_light.txt. You decide you don't need the Food\par
\noindent folder at all. Delete the entire Food folder and everything inside it.List everything left in your\par
\noindent Backpack to see what you have.\par

\noindent REFERENCES: Sl.No.                     Text Book/ Website Author        Canonical Ltd. Ubuntu Official\par
\noindent Publication 1                    als/command-line-for-be\par
\noindent https://ubuntu.com/tutori\par

\noindent ginners\par

\noindent 2                           htips:/www.geeksforgee      GeeksforGeeks            GeeksforGeeks\par
\noindent ks.org/linux-file-hierarch\par
\noindent y-structure/\par
\clearpage
\expheading{3}{\underline{\hspace{1.8cm}}}
\exptitle{WEB \& CONNECTIVITY}
\medskip\noindent\textbf{PRACTICAL SIGNIFICANCE}\par
\noindent To practice effective web searching using advanced search operators, understand email etiquette,\par
\noindent and familiarize students with the use of Google Classroom.\par
\noindent RELEVANT PROGRAM OUTCOMES(POS) PO1: Apply knowledge of mathematics, natural sciences,\par
\noindent engineering fundamentals, and an\par
\noindent engineering specialization (as specified in DK1 to DK4) to a wide range of practical procedures and\par
\noindent practices.\par

\medskip\noindent\textbf{RELEVANT COURSE OUTCOMES(COS)}\par
\noindent CO1L : Identify the basic functions and features of Computers, Operating Systems, and Internet\par
\noindent applications.\par

\medskip\noindent\textbf{THEORY}\par
\noindent Web connectivity refers to the ability of a computer system to connect to the internet and access\par
\noindent online resources and services.\par

\noindent Requirements for web connectivity:\par
\noindent\textbullet\; Computer or mobile device\par
\noindent\textbullet\; Internet connection\par
\noindent\textbullet\; Web browser\par
\noindent\textbullet\; User account (for online services)\par

\noindent Examples of web browsers:\par
\noindent\textbullet\; Google Chrome, Mozilla Firefox, Microsoft Edge, Safari\par
\noindent Advanced Search Operators\par
\noindent Search operators are special commands used with search engines to obtain more accurate and\par
\noindent filtered results.\par

\noindent Common search operators:\par
\noindent Operator Purpose Example "" Search exact phrase "computer networks" site: Search within a specific\par
\noindent website site:edu operating system filetype: Search for a particular file type filetype:pdf Linux tutorial\par
\noindent OR Search either keyword Python OR Java\par
\noindent - Exclude a word Java -coffee\par
\noindent\textbullet\; Wildcard operator "largest * in India” Advantages:\par
\noindent\textbullet\; Reduces irrelevant search results\par
\noindent\textbullet\; Saves time\par
\noindent\textbullet\; Provides accurate information\par
\noindent\textbullet\; Helps locate specific documents\par
\medskip\noindent\textbf{Email Etiquette}\par
\noindent Email etiquette refers to the professional and proper behavior followed while sending and receiving\par
\noindent emails.\par

\noindent Guidelines for proper email writing:\par
\noindent 1. Use a clear subject line Example:\par
\clearpage
\noindent Subject: Submission of Lab Record\par
\noindent 2. Begin with proper greeting\par
\noindent B\par

\noindent Examples:\par
\noindent\textbullet\; Dear Sir/Madam,\par
\noindent\textbullet\; Respected Sir,\par
\noindent\textbullet\; Hello Team,\par
\noindent 3. Write concise and clear messages\par
\noindent\textbullet\; Use simple language\par
\noindent\textbullet\; Keep the content relevant\par
\noindent\textbullet\; Avoid unnecessary information 4, Maintain professional language Avoid:\par
\noindent\textbullet\; Informal words\par
\noindent\textbullet\; Slang language\par
\noindent\textbullet\; Excessive punctuation\par
\noindent 5. End politely Examples:\par
\noindent\textbullet\; Thank you\par
\noindent\textbullet\; Regards\par
\noindent\textbullet\; Yours faithfully\par
\noindent 6. Check before sending\par
\noindent\textbullet\; Recipient address\par
\noindent\textbullet\; Attachments\par
\noindent\textbullet\; Grammar\par
\noindent\textbullet\; Spelling mistakes\par
\noindent Google Classroom\par
\noindent Google Classroom is an online learning platform used to manage classes and assignments digitally:\par

\noindent Features:\par
\noindent\textbullet\; Join classes using class code\par
\noindent\textbullet\; View assignments\par
\noindent\textbullet\; Submit assignments\par
\noindent\textbullet\; Download study materials\par
\noindent\textbullet\; Communicate with instructors\par
\noindent\textbullet\; Receive notifications Advantages:\par
\noindent\textbullet\; Easy content sharing\par
\noindent\textbullet\; Paperless submission\par
\noindent\textbullet\; Better communication\par
\noindent\textbullet\; Organized learning environment\par

\noindent RESOURCES REQUIRED SINo.               Computer System Desktop/Laptop    1 per student\par
\noindent [Name of resource Suggested            with Ubuntu\par
\noindent specification 1\par

\noindent Linux\par

\noindent 2                                                      Internet Connection Broadband / Wi-Fi 1 per lab\par

\noindent 3                                                      Web Browser Mozilla Firefox / Google Installed\par
\clearpage
\noindent Chrome\par

\noindent 4                                      Google Account Valid institution or   1 per student email ID\par
\noindent personal\par

\noindent 5                                                        Google Classroom Accessed via browser 1 per\par
\noindent student\par

\medskip\noindent\textbf{AI TEACHERS AND TRAINING}\par
\noindent PROCEDURE List the one-by-one procedure to be followed for completing the experiment. Activity 1:\par
\noindent Advanced Search Operators\par
\noindent 1. Open a web browser. 2, Open a search engine.\par
\noindent 3. Practice the following searches:\par
\noindent\textbullet\; "Computer Networks"\par
\noindent\textbullet\; site:edu computer engineering\par
\noindent\textbullet\; filetype:pdf Linux notes\par
\noindent\textbullet\; Python OR Java\par
\noindent\textbullet\; Java -coffee 4, Observe search results. Activity 2: Email Etiquette\par
\noindent 1. Open email application.\par
\noindent 2. Compose a new email.\par
\noindent 3. Enter recipient email address.\par
\noindent 4. Enter subject.\par
\noindent 5. Write a proper message.\par
\noindent 6. Attach a file if required.\par
\noindent 7. Verify details and send. Activity 3: Google Classroom\par
\noindent 1. Open Google Classroom.\par
\noindent 2. Login using Google account.\par
\noindent 3. Join a class using class code.\par
\noindent 4. Access learning materials,\par
\noindent 5. Submit a sample assignment.\par

\medskip\noindent\textbf{RESULTS}\par
\noindent Thus advanced search operators, email etiquette practices, and Google Classroom functionalities\par
\noindent were studied and practiced successfully.\par
\medskip\noindent\textbf{PRACTICAL RELA’}\par
\noindent 1. What is a search operator?\par
\noindent 2. What does site: operator do?\par
\noindent 3. What is filetype: operator?\par
\noindent 4, Why is email etiquette important?\par
\noindent 5. What is Google Classroom?\par
\noindent 6. Name two web browsers.\par

\medskip\noindent\textbf{SELF-LEARNING COMPONENTS}\par
\noindent 1. Use filetype:pdf "Computer Engineering” site:edu on Google. Record the URL and File Size (KB) of\par
\noindent the first PDF result found.\par
\noindent 2. Enroll 10 google classroom or any other LMS of Foundational IT Skills and customize profile.\par
\clearpage
\noindent 3. Create an email id, draft a formal mail and send mail to your friends with netiquitte.(scenario- You\par
\noindent are the Event Coordinator for your college. You need to send a ‘formal announcement to the Staff\par
\noindent coordinator, 10 student participants regarding a change in the workshop schedule. You must also\par
\noindent keep the HoD informed of this update, but you must ensure the students’ private email addresses\par
\noindent are not visible to one another:)\par
\medskip\noindent\textbf{REFERENCES,}\par
\noindent SL.No. ‘Text Book/ Website\par

\noindent 1                                              https://ahrefs.comvblog/google-advanced-search-oper\par
\noindent ators!\par
\clearpage
\expheading{4}{\underline{\hspace{1.8cm}}}
\exptitle{PROFESSIONAL DOCUMENTATION}
\medskip\noindent\textbf{PROFESSIONAI DOCUMENTATION}\par

\medskip\noindent\textbf{PRACTICAL SIGNIFICANCE}\par
\noindent To create a professional document using a word processing application with proper ‘formatting,\par
\noindent styles, and an automatically generated table of contents.\par

\medskip\noindent\textbf{RELEVANT PROGRAM OUTCOMES(POS)}\par
\noindent PO1: Apply knowledge of mathematics, natural sciences, engineering fundamentals, and an\par
\noindent engineering specialization (as specified in DK1 to DK4) to a wide range of practical procedures and\par
\noindent practices.\par
\noindent PO6: When solving well-defined engineering problems, evaluate sustainable development\par
\noindent impacts to: society, the economy, sustainability, health and safety, legal frameworks, and the\par
\noindent environment (DK1, DK5, and DK7).\par
\noindent POT: Understand and commit to professional ethics and the norms of technician practice, including\par
\noindent compliance with relevant laws. Demonstrate an understanding of the need for diversity and inclusion\par
\medskip\noindent\textbf{(DK9).}\par

\medskip\noindent\textbf{RELEVANT COURSE OUTCOMES(COS)}\par
\noindent CO2: Use Stand-alone and Cloud-based office tools to prepare documents, spreadsheets, and\par
\noindent presentations.\par

\medskip\noindent\textbf{THEORY}\par
\noindent Professional documentation refers to preparing documents in a structured and organized ‘manner for\par
\noindent communication and presentation purposes.\par

\noindent Examples:\par
\noindent\textbullet\; Project reports, Seminar reports, Lab records, Technical documentation, Business reports\par
\noindent Characteristics of a professional document:\par
\noindent\textbullet\; Proper title and headings\par
\noindent\textbullet\; Consistent formatting\par
\noindent\textbullet\; Page numbering\par
\noindent\textbullet\; Table of contents\par
\noindent\textbullet\; Proper font and alignment\par
\noindent\textbullet\; Well-organized sections\par
\noindent Word Processing Software\par
\noindent Word processing software allows users to create, edit, format, and print documents, Examples:\par
\noindent\textbullet\; Microsoft Word\par
\noindent\textbullet\; LibreOffice Writer\par
\noindent\textbullet\; Google Docs\par

\noindent Functions:\par
\noindent\textbullet\; Text editing\par
\noindent\textbullet\; Formatting\par
\noindent\textbullet\; Inserting images and tables\par
\clearpage
\noindent\textbullet\; Spell checking\par
\noindent\textbullet\; Document organization\par

\noindent Styles\par
\noindent Styles are predefined formatting settings used to maintain consistency throughout a document.\par

\noindent Common styles:\par
\noindent Style Purpose\par
\noindent Title Main document title Heading 1 Main headings Heading 2 Subheadings\par
\noindent Normal\par
\noindent Regular text\par
\noindent Subtitle\par
\noindent Secondary titles\par

\noindent Advantages of using styles:\par
\noindent\textbullet\; Uniform appearance\par
\noindent\textbullet\; Easy formatting\par
\noindent\textbullet\; Faster editing\par
\noindent\textbullet\; Automatic table of contents generation\par

\noindent Table of Contents (TOC)\par
\noindent A table of contents is a list of document headings and corresponding page numbers.\par

\noindent Advantages:\par
\noindent\textbullet\; Helps locate sections quickly\par
\noindent\textbullet\; Improves document organization\par
\noindent\textbullet\; Provides professional appearance\par
\noindent\textbullet\; Updates automatically when document content changes\par

\medskip\noindent\textbf{RESOURCES REQUIRED}\par
\noindent Sl.No.                       \_ Name of resource        Computer System            1 per student\par
\noindent Suggested specification   Desktop/Laptop with\par
\noindent Quantity 1                Ubuntu\par

\noindent Linux 2 LibreOffice          Google Account Valid      1 per student email ID 5   Google Classroom\par
\noindent Writer Version 7.0 or        institution or personal                              Accessed via browser 1\par
\noindent above Installed 3 Internet                                                        per student\par
\noindent Connection Broadband /\par
\noindent Wi-Fi 1 per lab 4\par

\noindent AIN\par
\medskip\noindent\textbf{SCHEME}\par

\medskip\noindent\textbf{PROCEDURE}\par
\noindent Activity: Creating a Formal Report\par
\noindent 1. Open the word processing application.\par
\noindent 2. Create a new document.\par
\noindent 3. Enter report ttle. 4, Add headings and content.\par
\noindent 5. Apply styles:\par
\noindent\textbullet\; Title style for title\par
\noindent\textbullet\; Heading 1 for main headings\par
\clearpage
\noindent\textbullet\; Heading 2 for subheadings\par
\noindent\textbullet\; Normal style for body text\par
\noindent 6. Insert page numbers.\par
\noindent 7. Generate an automatic table of contents.\par
\noindent 8. Save the document.\par
\noindent Sample Report Structure Title:Computer Fundamentals Report\par
\noindent 1. Introduction\par
\noindent 2. Hardware Components\par
\noindent 3. Software Components 4, Advantages of Computers\par
\noindent 5. Applications of Computers\par
\noindent 6. Conclusion\par
\noindent RE!\par
\noindent He\par
\noindent Thus a formal report was created successfully using styles and an automatically generated table of\par
\noindent contents\par
\medskip\noindent\textbf{PRACTICAL RELATED QUESTPRACTICAL RELATED QUESTPRACTICAL RELATED QUESTIONS 1,}\par
\noindent What is professional documentation?\par
\noindent 2. What is a style?\par
\noindent 3. Why are styles used?\par
\noindent 4, What is TOC?\par
\noindent 5. Which style is used for main headings?\par
\noindent 6. Name two word processing software applications.\par
\noindent 7. Why is formatting important?\par
\noindent 8. How can a table of contents be generated?\par
\noindent 9. What is the use of page numbering?\par
\noindent 10, State one advantage of professional documentation.\par
\medskip\noindent\textbf{SELF- LEARNING COMPONENTS,}\par
\noindent 1. Create a standard I-page "Assignment Cover Sheet" using Word Styles (Heading 1, Heading 2).\par
\noindent 2. Create a professional Biodata for an internship application and export it into pdf (The\par
\noindent Specification:\par
\noindent Page Layout: Margins: Set to "Normal" (1 inch or 2.54 cm on all sides). Font: Use a professional\par
\noindent sans-serif font (e.g., Arial or Calibri) for the body.\par
\noindent Use of Styles: Apply Heading 1 10 main sections (Personal Details, Education, Skills). Modify the\par
\noindent Heading 1 style to be Bold and Underlined,\par
\noindent Data Organization (The Table): Use a 2-column Table to align your personal information (e.g., Date of\par
\noindent Birth, Nationality, Languages) so the colons (:) line up perfectly:\par
\noindent Border: Remove the table borders so the text looks floating but perfectly aligned. Visual\par
\noindent Element:Insert a placeholder for a professional photo.\par
\medskip\noindent\textbf{‘REFERENCES}\par
\noindent SIN. Text Book/ Website ‘Author Publication\par

\noindent 1                 https://help libr   The Document   LibreOffice n-U    Foundation   Official Help\par
\noindent eoffice.org/lates                  S/text/swriter/g\par
\noindent ve                                 uide/main.html\par
\clearpage
\null
\thispagestyle{main}
\clearpage
\clearpage
\expheading{5}{\underline{\hspace{1.8cm}}}
\exptitle{SPREADSHEET BASICS}
\medskip\noindent\textbf{PRACTICAL SIGNIFICANCE}\par
\noindent To understand spreadsheet operations and use formulas such as SUM, AVERAGE (AVG), COUNT,\par
\noindent COUNTIF, and IF with cell referencing.\par

\medskip\noindent\textbf{RELEVANT PROGRAM OUTCOMES(POS)}\par
\noindent PO1: Apply knowledge of mathematics, natural sciences, engineering fundamentals, and an\par
\noindent engineering specialization (as specified in DK1 to DK4) to a wide range of practical procedures and\par
\noindent practices.\par
\noindent PO6:; When solving well-defined engineering problems, evaluate sustainable development impacts to:\par
\noindent society, the economy, sustainability, health and safety, legal frameworks, and the environment (DK1,\par
\noindent DK5, and DK7).\par
\noindent PO7: Understand and commit to professional ethics and the norms of technician practice, including\par
\noindent compliance with relevant laws. Demonstrate an understanding of the need for diversity and inclusion\par
\medskip\noindent\textbf{(DK9).}\par

\medskip\noindent\textbf{RELEVANT COURSE OUTCOMES(COS)}\par
\noindent CO2: Use Stand-alone and Cloud-based office tools to prepare documents, spreadsheets, and\par
\noindent presentations\par

\medskip\noindent\textbf{THEORY}\par
\noindent Spreadsheet\par
\noindent A spreadsheet is an electronic worksheet used for organizing, storing, and performing calculations on\par
\noindent data.\par
\noindent Examples of spreadsheet software: Microsoft Excel, LibreOffice Calc, Google Sheets Applications of\par
\noindent spreadsheets:\par
\noindent Mark list preparation\par
\noindent Attendance records\par
\noindent\textbullet\; Payroll processing\par
\noindent Budget preparation\par
\noindent Data analysis\par
\medskip\noindent\textbf{(TEACHERS AND ERY ~ DIPLOMA 2}\par
\noindent ME\par

\noindent Components of a Spreadsheet\par
\noindent Workbook\par
\noindent A workbook is a collection of worksheets.\par

\noindent Worksheet\par
\noindent A worksheet consists of rows and columns used for entering data.\par

\noindent Row\par
\noindent Horizontal arrangement of cells\par
\noindent Identified by numbers\par
\clearpage
\noindent Example:\par
\noindent 12,3...\par

\noindent Column\par
\noindent Vertical arrangement of cells\par
\noindent Identified by letters\par

\noindent Example:\par
\medskip\noindent\textbf{ABC,}\par

\noindent Cell\par
\noindent Intersection of row and column\par
\noindent Stores data\par

\noindent Example:\par
\medskip\noindent\textbf{AI, B2, C5}\par

\noindent Cell Referencing\par
\noindent Cell referencing is used to refer to a cell address in formulas. Types of cell references:\par
\noindent 1. Relative Reference\par

\noindent Changes automatically when copied. Example:\par
\medskip\noindent\textbf{=AI+B1}\par
\noindent 2. Absolute Reference\par
\noindent Does not change when copied.\par

\noindent Uses \$ symbol. Example:\par
\medskip\noindent\textbf{=SASI+SBS1}\par
\noindent 3. Mixed Reference\par

\noindent Either row or column remains fixed. Examples:\par
\medskip\noindent\textbf{=SAI}\par
\medskip\noindent\textbf{=ASI}\par

\noindent Spreadsheet Functions\par
\noindent SUM Function\par
\noindent Used to calculate total values. Syntax:\par
\noindent =SUM(range)\par

\noindent Example:\par
\medskip\noindent\textbf{=SUM(B2:B6)}\par

\noindent AVERAGE Function (AVG)\par
\noindent Used to calculate average values. Syntax:\par
\noindent AVERAGE (range)\par

\noindent Example:\par
\medskip\noindent\textbf{=AVERAGE(B2:B6)}\par

\noindent COUNT Function\par
\noindent Counts the number of cells containing numeric values Syntax:\par
\noindent =COUNT(range)\par
\clearpage
\noindent Example:\par
\medskip\noindent\textbf{=COUNT(B2:B6)}\par

\noindent COUNTIF Function\par
\noindent Counts cells that satisfy a condition.\par
\noindent Syntax: =COUNTIF (range,condition)\par

\noindent Example:\par
\medskip\noindent\textbf{=COUNTIF(B2:B6,">50")}\par

\noindent IF Function\par
\noindent Performs logical testing.\par

\noindent Syntax:\par
\noindent =IF (condition, value\_if.true,value\_if false) Example: IF (B2>\par
\noindent (0,"Pass”,"Fail”)\par

\noindent SI.No. [ Name of            Suggested specification   Computer System          Desktop/Laptop with 1\par
\noindent resource                    Quantity 1                                         per student\par

\noindent Ubuntu Linux 2        LibreOffice Cale/     LibreOffice Cale     Google Classroom   Accessed via\par
\noindent Installed Google                        browser 1 per\par
\noindent sheets Version 7.0                      student\par
\noindent or above 3\par

\medskip\noindent\textbf{PROCEDURE}\par
\noindent Activity: Using Spreadsheet Functions\par
\noindent Open spreadsheet software. Create a new worksheet.\par
\noindent Enter student details and marks. Apply formulas.\par
\noindent SUM\par
\medskip\noindent\textbf{AVERAGE}\par
\medskip\noindent\textbf{COUNT}\par
\medskip\noindent\textbf{COUNTIF}\par
\noindent IF\par
\noindent Apply cell references.\par
\noindent Observe the results.\par
\noindent Save the worksheet.\par
\noindent NAMNeeeee Kune\par
\medskip\noindent\textbf{(TEACHERS AND ERY ~ DIPLOMA 2}\par
\medskip\noindent\textbf{ME *}\par

\medskip\noindent\textbf{RESULTS}\par
\noindent Spreadsheet formulas (SUM, AVERAGE, COUNT, COUNTIF, and IF) and cell referencing\par
\noindent were studied and applied successfully and uploaded the spreadsheet in the Google classroom.\par
\medskip\noindent\textbf{PRACTICAL RELATED QUESTPRACTICAL RELATED QUESTPRACTICAL RELATED QUESTIONS}\par
\noindent What is a spreadsheet?\par
\clearpage
\noindent What is a workbook?\par
\noindent What is a worksheet?\par
\noindent What is cell referencing?\par
\noindent What is the syntax of SUM function? What is COUNT function?\par
\noindent What is COUNTIF used for?\par
\noindent What is an IF function?\par
\noindent 10. Name two spreadsheet applications.\par

\medskip\noindent\textbf{SELF-LEARNING COMPONENTS}\par
\noindent 1. Find and apply the COUNTIF formula in a spreadsheet to count how many students\par
\noindent passed a test\par
\noindent 2. In a cricket database, use the SUM and AVERAGE functions to calculate the total and mean values\par
\noindent for the "Runs Scored" field in the range B2:B21. Next, apply the COUNTIF\par
\noindent Differentiate between relative and absolute references.\par
\noindent formula to the "Player Category" field in range C2:C21 to find the total number of players listed as\par
\noindent "Batsman.” In the "Selection Status" column, use an IF formula to check the "Runs\par
\noindent Scored” in cell B2 and display "Selected" if the value is 50 or higher, otherwise display "Waitlist."\par
\noindent Finally, use the COUNT function on the "Player Name" field in range A2:A21 to\par
\noindent verify the total number of participants in the academy.\par

\noindent REFEREN     Text         https://hel   The        LibreOffic     Foundatio    hups:/ww    Geeksfor   Geeksfor\par
\noindent CES         Book/        p.libreoffi   Document   e Official l   n Help de/   w.geeksfo   Geeks      Geeks\par
\noindent Sl.No.      Website      ce.org/                  atesten-U      main.html    rgeeks\par
\noindent ‘Author P                             S/text/sca     2\par
\noindent ublication                            le/gui\par
\noindent 1\par

\noindent org/basic-excel-formulas/\par
\noindent a\par
\clearpage
\expheading{6}{\underline{\hspace{1.8cm}}}
\exptitle{PRESENTATION DESIGN}
\medskip\noindent\textbf{PRACTICAL SIGNIFICANCE}\par
\noindent To create a professional 5-slide pitch deck using LibreOffice Impress in Ubuntu by applying themes,\par
\noindent animations, and slide transitions.\par

\medskip\noindent\textbf{RELEVANT PROGRAM OUTCOMES(POS)}\par
\noindent PO1: Apply knowledge of mathematics, natural sciences, engineering fundamentals, and an\par
\noindent engineering specialization (as specified in DK1 to DK4) to a wide range of practical procedures and\par
\noindent practices.\par
\noindent PO6: When solving well-defined engineering problems, evaluate sustainable development impacts to:\par
\noindent society, the economy, sustainability, health and safety, legal frameworks, and the environment (DK1,\par
\noindent DK5, and DK7).\par
\noindent PO7: Understand and commit to professional ethics and the norms of technician practice, including\par
\noindent compliance with relevant laws. Demonstrate an understanding of the need for diversity and inclusion\par
\medskip\noindent\textbf{(DK9).}\par

\medskip\noindent\textbf{RELEVANT COURSE OUTCOMES(COS)}\par
\noindent CO2: Use Stand-alone and Cloud-based office tools to prepare documents, spreadsheets, and\par
\noindent presentations.\par
\medskip\noindent\textbf{‘THEORY}\par

\noindent Presentation Software\par
\noindent Presentation software is used to create visual presentations containing text, images, graphics, charts,\par
\noindent and multimedia objects.\par

\noindent Examples:\par
\noindent © LibreOffice Impress\par
\noindent © Microsoft PowerPoint\par
\noindent © Google Slides\par

\noindent Pitch Deck\par
\noindent A pitch deck is a short presentation used to communicate ideas, projects, products, or solutions\par
\noindent effectively.\par
\noindent A typical pitch deck contains\par
\noindent Title\par
\noindent Problem Statement\par
\noindent Solution\par
\noindent Features/Benefits\par
\noindent Conclusion\par
\noindent Animation Animation is the visual effect applied to objects within a slide.\par

\noindent Examples:\par
\noindent Appear\par
\clearpage
\noindent Fade In\par
\noindent Flyin\par
\noindent Zoom\par

\noindent Advantages:\par
\noindent © Makes presentations attractive\par
\noindent © Highlights important information\par
\noindent © Improves audience engagement\par
\noindent Slide Transition Slide transition is the visual effect displayed when moving from one slide to another.\par
\noindent Examples:\par
\noindent Fade\par
\noindent Push\par
\noindent Wipe\par
\noindent Split\par

\noindent Advantages:\par
\noindent © Improves presentation flow\par
\noindent © Gives a professional appearance\par

\medskip\noindent\textbf{RESOURCES REQUIRED}\par
\noindent Sl.No.                      Name of resource          Computer System            1 per student\par
\noindent Suggested specification   Desktop/Laptop with\par
\noindent Quantity 1                Ubuntu\par

\noindent Linux 2                                               LibreOffice Impress Version 7.0 or above Installed\par

\medskip\noindent\textbf{PROCEDURE}\par
\noindent Activity: Creating a 5-Slide Pitch Deck\par
\noindent Open LibreOffice Impress. Create a new presentation. Select a suitable theme. Create five slides.\par
\noindent Add content to each slide. Insert images if required. Apply animations to objects. Apply slide\par
\noindent transitions.\par
\noindent 1. Save the presentation.\par
\noindent 10. Run the slideshow.\par
\noindent SEN AwAWNS\par

\medskip\noindent\textbf{RESULTS}\par
\noindent A visual presentation containing text, images, graphics and multimedia objects was created\par
\noindent successfully using presentation software and uploaded the pdf format of the same in the Google\par
\noindent Classroom.\par
\medskip\noindent\textbf{PRACTICAL RELATED QUESTPRACTICAL RELATED QUESTPRACTICAL RELATED QUESTIONS}\par
\noindent What is presentation software?\par
\noindent What is a pitch deck?\par
\noindent What is animation?\par
\noindent What is a transition?\par
\noindent Name two animation effects.\par
\clearpage
\noindent ‘Name two transition effects.\par
\noindent Which application is used in Ubuntu for presentations? What is the file extension of LibreOffice\par
\noindent Impress? Which key starts slideshow mode?\par
\noindent 10. State one advantage of presentation software.\par
\medskip\noindent\textbf{SEN AWAWNS}\par

\medskip\noindent\textbf{SELF-LEARNING COMPONENTS}\par
\noindent 1, Edit a "Slide Master" to ensure your name and roll number appear automatically on header or\par
\noindent footer of every slide of a deck.\par
\noindent REFERENCES SL No.\par

\noindent 1                  hutps:/’help.libr   The Document   \_LibreOftice sen   Foundation   Official Help\par
\noindent eoftice.org/late                   -US/text/simpre                 main.html\par
\noindent ss/guide/\par

\noindent Text Book/ Website Author Publication\par
\clearpage
\expheading{7}{\underline{\hspace{1.8cm}}}
\exptitle{CLOUD COLLABORATION}
\medskip\noindent\textbf{PRACTICAL SIGNIFICANCE}\par
\noindent To understand cloud-based collaboration by using real-time editing features in Google Docs and\par
\noindent Google Sheets and to explore version history management.\par

\medskip\noindent\textbf{RELEVANT PROGRAM OUTCOMES(POS)}\par
\noindent PO1: Apply knowledge of mathematics, natural sciences, engineering fundamentals, and an\par
\noindent engineering specialization (as specified in DK1 to DK4) to a wide range of practical procedures and\par
\noindent practices.\par
\noindent PO6: When solving well-defined engineering problems, evaluate sustainable development impacts to:\par
\noindent society, the economy, sustainability, health and safety, legal frameworks, and the environment (DK1,\par
\noindent DK5, and DK7)\par
\noindent POT: Understand and commit to professional ethics and the norms of technician practice, including\par
\noindent compliance with relevant laws. Demonstrate an understanding of the need for diversity and inclusion\par
\medskip\noindent\textbf{(DK9).}\par

\medskip\noindent\textbf{RELEVANT COURSE OUTCOMES(COS)}\par
\noindent CO2: Use Stand-alone and Cloud-based office tools to prepare documents, spreadsheets, and\par
\noindent presentations\par
\noindent THEORY ‘Cloud Collaboration\par
\noindent Cloud collaboration is the process of allowing multiple users to access, edit, and share files or\par
\noindent resources through internet-based platforms.\par

\noindent ‘Advantages:\par
\noindent\textbullet\; Enables real-time teamwork\par
\noindent\textbullet\; Automatic saving of work\par
\noindent\textbullet\; Access from any device\par
\noindent\textbullet\; Easy sharing of files\par
\noindent\textbullet\; Maintains version history\par

\noindent Examples of cloud collaboration platforms:\par
\noindent\textbullet\; Google Docs\par
\noindent\textbullet\; Google Sheets\par
\noindent\textbullet\; Google Drive\par
\noindent\textbullet\; Microsoft OneDrive\par
\noindent\textbullet\; Dropbox\par

\noindent Google Docs\par
\noindent Google Docs is a web-based document editor used to create and edit text documents online. Features:\par
\noindent\textbullet\; Real-time editing\par
\noindent\textbullet\; Auto-save\par
\noindent\textbullet\; Sharing and permissions\par
\noindent\textbullet\; Comments and suggestions\par
\noindent\textbullet\; Version history\par

\noindent Applications:\par
\clearpage
\noindent\textbullet\; Report writing\par
\noindent\textbullet\; Documentation\par
\noindent\textbullet\; Group assignments\par
\noindent\textbullet\; Project work\par

\medskip\noindent\textbf{Google Sheets}\par
\noindent Google Sheets is an online spreadsheet application used to create and analyze data. Features:\par
\noindent\textbullet\; Spreadsheet calculations\par
\noindent\textbullet\; Real-time collaboration\par
\noindent\textbullet\; Formula support\par
\noindent\textbullet\; Chart creation\par
\noindent\textbullet\; Version history\par

\noindent Applications:\par
\noindent\textbullet\; Mark lists\par
\noindent\textbullet\; Attendance sheets\par
\noindent\textbullet\; Budget preparation\par
\noindent\textbullet\; Data analysis\par

\noindent Version History\par
\noindent Version history stores previous changes made to a document. Advantages:\par
\noindent\textbullet\; Track changes made by users\par
\noindent\textbullet\; Restore previous versions\par
\noindent\textbullet\; Identify contributors\par
\noindent\textbullet\; Prevent accidental data loss\par

\medskip\noindent\textbf{RESOURCES REQUIRED}\par
\noindent SI. No. Name of resource Suggested specification Quantity\par

\noindent 1                                      Computer System Desktop/Laptop       1 per student\par
\noindent with Ubuntu\par

\noindent Linux 2         internet        Web Browser      \_ Installed 4   Google           1 per student   Google Docs\par
\noindent Connection      Mozilla                          Account          1D 5            Accessed via\par
\noindent Broadband /     Firefox /                        ‘Valid                           browser 1\par
\noindent Wi-Fi 1 per     Google                           institution or                   per student\par
\noindent lab 3           Chrome                           personal                         (docs\par
\noindent email                            google.com)\par

\noindent 6                                      Google Sheets Accessed via           Google Classroom Accessed via\par
\noindent browser 1 per student (sheets        browser 1 per student\par
\noindent google.com) 7\par

\medskip\noindent\textbf{PROCEDURE}\par
\noindent Activity: Real-Time Editing and Version History\par
\noindent 1. Open a web browser in Ubuntu.\par
\noindent 2. Login to Google account.\par
\noindent 3. Open Google Docs.\par
\noindent 4. Create a new document.\par
\noindent 5. Enter sample content.\par
\noindent 6. Share the document with another user.\par
\clearpage
\noindent 7. Perform collaborative editing.\par
\noindent 8, Add comments and suggestions.\par
\noindent 9. Open Google Sheets and create sample data. 10, View version history.\par
\noindent II. Restore an earlier version if required.\par

\medskip\noindent\textbf{RESULTS}\par
\noindent Thus real-time collaboration using Google Docs and Google Sheets along with version history features\par
\noindent was studied and practiced successfully and uploaded the report in Google classroom,\par
\medskip\noindent\textbf{PRACTICAL RELATED QUESTPRACTICAL RELATED QUESTPRACTICAL RELATED QUESTIONS}\par
\noindent 1. What is cloud collaboration?\par
\noindent 2. What is Google Docs?\par
\noindent 3. What is Google Sheets?\par
\noindent 4. What is version history?\par
\noindent 1008 - Foundational IT\par
\noindent Skills Semester 1\par
\noindent 5. What is real-time editing?\par
\medskip\noindent\textbf{SELF- LEARNING}\par
\medskip\noindent\textbf{COMPONENTS}\par
\noindent 1. Imagine a team of four students is editing a single Google Sheet for a lab report. One student\par
\noindent accidentally deletes a large data table and then closes the browser. Explain how the team can use\par
\noindent the ‘Version History’ tool to see the names of people who edited the file and describe the steps to\par
\noindent restore the table back to its original state.\par
\noindent 2, In Google Sheets, if three people are working on a project, how can you tell which specific person\par
\noindent deleted a value in Cell BS? Name the menu path you would use to find this information and explain\par
\noindent how the ‘color-coded’ system helps identify the person.\par
\noindent 3. Create a I-page overleaf document that includes a title, a mathematical formula, and a\par
\noindent list.(Optional Assignment)\par
\noindent REFERENCES No. Text Book/ Website ‘Author [Publication\par

\noindent 1 htips://support.google.com/          Google                    Google Docs does/answer/2494822\par
\noindent Help Centre\par

\noindent 2                           https://support.google.co   Google           Google Docs\par
\noindent m/                                           docs/answer/190843\par
\noindent Help Centre\par
\clearpage
\expheading{8}{\underline{\hspace{1.8cm}}}
\exptitle{DATA ANALYSIS}
\medskip\noindent\textbf{DATA ANAIYSIS}\par

\medskip\noindent\textbf{PRACTICAL SIGNIFICANCE}\par
\noindent To analyze a dataset using Pivot Tables and Slicers in Google Sheets for summarizing and \{filtering\par
\noindent information.\par
\medskip\noindent\textbf{RELEVANT PROGRAM OUTCOMES (POS)}\par
\noindent PO1: Apply knowledge of mathematics, natural sciences, engineering fundamentals, and an\par
\noindent engineering specialization (as specified in DK1 to DK4) to a wide range of practical procedures and\par
\noindent practices\par
\noindent PO2: Identify and analyze well-defined engineering problems and reach substantiated conclusions\par
\noindent using codified methods of analysis specific to their field of activity (DK1 to DK4).\par
\noindent PO4: Conduct investigations of well-defined problems by locating and search relevant codes and\par
\noindent catalogues, and by conducting standard tests and measurements (DK8).\par
\noindent PO7: Understand and commit to professional ethics and the norms of technician practice, including\par
\noindent compliance with relevant laws. Demonstrate an understanding of the need for diversity and inclusion\par
\medskip\noindent\textbf{(DK9)}\par

\medskip\noindent\textbf{RELEVANT COURSE OUTCOMES(COS)}\par
\noindent CO3: Apply spreadsheet functions and Generative AI tools to construct interactive dashboards,\par
\noindent automated reports, and professional visual presentations for interpreting complex datasets.\par

\medskip\noindent\textbf{THEORY}\par
\noindent Data analysis is the process of organizing and interpreting data to extract useful information,\par
\noindent Applications.\par
\noindent\textbullet\; Student performance analysis\par
\noindent\textbullet\; Attendance analysis,\par
\noindent\textbullet\; Sales analysis\par
\noindent\textbullet\; Survey reports\par
\noindent\textbullet\; Budget planning\par

\noindent Pivot Table\par
\noindent A Pivot Table is a spreadsheet feature used to summarize and rearrange large amounts of data\par
\noindent dynamically.\par

\noindent Functions:\par
\noindent\textbullet\; Calculate totals\par
\noindent\textbullet\; Calculate averages\par
\noindent\textbullet\; Count records\par
\noindent\textbullet\; Group data\par
\noindent\textbullet\; Rearrange information Advantages:\par
\noindent\textbullet\; Faster analysis\par
\noindent\textbullet\; Easy report generation\par
\noindent\textbullet\; Better understanding of data\par

\noindent Slicer\par
\clearpage
\noindent A Slicer is a visual filtering tool used with Pivot Tables. Functions:\par
\noindent\textbullet\; Filters data interactively\par
\noindent\textbullet\; Displays selected categories\par
\noindent\textbullet\; Makes data analysis easier Advantages:\par
\noindent\textbullet\; User-friendly\par
\noindent\textbullet\; Fast filtering\par
\noindent\textbullet\; Improves visualization\par

\medskip\noindent\textbf{RESOURCES REQUIRED}\par
\noindent SINo. Name of resource Suggested specification Quantity\par

\noindent 1                                      Computer System Desktop/Laptop    1 per student\par
\noindent with Ubuntu\par

\noindent Linux\par

\noindent 2                                                      Internet Connection Broadband / Wi-Fi 1 per lab\par

\noindent 3                                                      Web Browser Mozilla Firefox / Google Chrome \} \_\par
\noindent Installed\par

\noindent 4                                                      Google Account Valid institution or personal 1 per\par
\noindent student\par

\noindent email ID 5                                             Google Sheets Accessed via browser 1 per student\par
\noindent (sheets. google.com)\par

\noindent 6                                                      Google Classroom Accessed via browser 1 per\par
\noindent student\par

\medskip\noindent\textbf{PROCEDURE}\par
\noindent Activity: Data Analysis Using Pivot Tables and Slicers\par
\noindent 1. Open Google Sheets.\par
\noindent 2. Create a blank spreadsheet.\par
\noindent 3. Enter sample data.\par
\noindent 4, Select the dataset.\par
\noindent 5. Create a Pivot Table.\par
\noindent 6. Configure row, column, and value fields.\par
\noindent 7, Adda Slicer.\par
\noindent 8, Filter the data.\par
\noindent 9, Analyze the results.\par
\noindent 10. Save the spreadsheet.\par
\noindent The dataset was successfully analyzed using Pivot Tables and Slicers in Google Sheets to summarize,\par
\noindent organize, and filter information. The observations and analysis results were\par
\medskip\noindent\textbf{(TEACHERS AND ERY ~ DIPLOMA 2}\par
\noindent ME\par
\noindent documented in a report containing the dataset, Pivot Table output, Slicer-based filtering results,\par
\noindent interpretations, and necessary screenshots. The completed report was successfully uploaded and\par
\noindent submitted through Google Classroom.\par
\clearpage
\medskip\noindent\textbf{PRACTICAL RELATED QUESTPRACTICAL RELATED QUESTIONS}\par
\noindent 1. What is a Pivot Table?\par
\noindent 2. What is a Slicer?\par
\noindent 3. Why is data analysis performed?\par
\noindent 4. What is filtering?\par
\noindent 5. What is the purpose of Pivot Tables?\par
\noindent 6. Name one advantage of Slicers.\par
\noindent 7. What is the use of row fields?\par
\noindent 8. What is the use of value fields?\par
\noindent 9. Which application was used for creating slicers and pivot table?\par
\noindent 10. State one application of data analysis.\par

\medskip\noindent\textbf{SELF-LEARNING COMPONENTS}\par
\noindent 1. Create a Pivot Table from a 50-row sales dataset. Drag "Category" to Rows and "Sales" to Values.\par
\noindent Record the Grand Total value displayed.\par
\noindent 2. You have a list of sales for ‘Laptop’, ‘Mouse’, and ‘Keyboard’. Create a Pivot Table to show the total\par
\noindent sales for each. Insert a Slicer for ‘Region’. If you click ‘South’ on the Slicer, what happens to the\par
\noindent ‘Grand Total’? If you add a new row of data to your original list, name the button you must click to\par
\noindent update the Pivot Table.\par
\medskip\noindent\textbf{REFERENCES,}\par
\noindent Sl.No. [Text Book/ Website ‘Author Publication\par

\noindent 1                           https://support.google.co   Google —\_        Google Sheets\par
\noindent mvdoes/answer/1272900\par
\noindent Help Centre\par

\noindent 2                           hups:/www.geeksforgee       GeeksforGeeks    GeeksforGeeks\par
\noindent ks.org/google-sheets/ho\par
\noindent ww-to-create-and-use-piv\par
\noindent o\par
\noindent t-tables-in-google-sheets/\par
\clearpage
\expheading{9}{\underline{\hspace{1.8cm}}}
\exptitle{DATA VISUALIZATION}
\medskip\noindent\textbf{PRACTICAL SIGNIFICANCE}\par
\noindent To create and analyze different types of charts in a spreadsheet for graphical representation and\par
\noindent interpretation of data.\par

\medskip\noindent\textbf{RELEVANT PROGRAM OUTCOMES(POS)}\par
\noindent PO1: Apply knowledge of mathematics, natural sciences, engineering fundamentals, and an\par
\noindent engineering specialization (as specified in DK1 to DK4) to a wide range of practical procedures and\par
\noindent practices.\par
\noindent PO2: Identify and analyze well-defined engineering problems and reach substantiated conclusions\par
\noindent using codified methods of analysis specific to their field of activity (DK1 to DK4).\par
\noindent PO4: Conduct investigations of well-defined problems by locating and search relevant codes and\par
\noindent catalogues, and by conducting standard tests and measurements (DK8).\par
\noindent PO7: Understand and commit to professional ethics and the norms of technician practice, including\par
\noindent compliance with relevant laws. Demonstrate an understanding of the need for diversity and inclusion\par
\medskip\noindent\textbf{(DK9).}\par

\medskip\noindent\textbf{RELEVANT COURSE OUTCOMES(COS)}\par
\noindent CO3: Apply spreadsheet functions and Generative AI tools to construct interactive dashboards,\par
\noindent automated reports, and professional visual presentations for interpreting complex datasets.\par

\medskip\noindent\textbf{THEORY}\par
\noindent Write the minimum theory required for performing the experiment.\par
\noindent Data visualization is the graphical representation of data and information using charts, graphs, and\par
\noindent diagrams to understand trends, patterns, and comparisons.\par

\noindent Advantages:\par
\noindent\textbullet\; Simplifies large amounts of data\par
\noindent\textbullet\; Helps identify trends and patterns\par
\noindent\textbullet\; Makes comparisons easier\par
\noindent\textbullet\; Improves understanding\par
\noindent\textbullet\; Supports decision-making Applications:\par
\noindent\textbullet\; Student performance analysis\par
\noindent\textbullet\; Sales analysis\par
\noindent\textbullet\; Attendance analysis\par
\noindent\textbullet\; Survey reports\par
\noindent\textbullet\; Financial reports\par

\noindent Types of Charts\par
\noindent 1. Column Chart Displays data using vertical bars. Applications:\par
\noindent\textbullet\; Comparing values among categories\par
\noindent\textbullet\; Comparing student marks\par
\noindent 2. Bar Chart Displays data using horizontal bars. Applications:\par
\noindent\textbullet\; Ranking information\par
\noindent\textbullet\; Comparing categories with longer names\par
\clearpage
\noindent 3. Pie Chart Displays percentage distribution of values. Applications:\par
\noindent\textbullet\; Showing proportions\par
\noindent\textbullet\; Category contribution analysis\par
\medskip\noindent\textbf{4. Line Chart}\par

\noindent Displays trends and changes over time. Applications:\par
\noindent\textbullet\; Attendance trends\par
\noindent\textbullet\; Monthly performance analysis\par
\noindent\textbullet\; Sales reports\par
\noindent 5. Seatter Chart\par

\noindent Displays relationships between two variables. Applications:\par
\noindent\textbullet\; Correlation studies\par
\noindent\textbullet\; Data comparison\par

\medskip\noindent\textbf{RESOURCES REQUIRED}\par
\noindent SI. No. ‘Name of resource Suggested specification       Computer System Desktop/Laptop with Ubuntu\} 1\par
\noindent Quantity 1                                              per student\par

\noindent Linux 2         Web Browser     \_ Installed 4   Google           1 per student   Google         Google\par
\noindent [Internet       Mozilla                         Account Valid    email ID 5      Sheets         Classroom\par
\noindent Connection      Firefox /                       institution or                   Accessed via   Accessed via\par
\noindent Broadband /     Google                          personal                         browser 1      browser 1\par
\noindent Wi-Fi 1 per     Chrome                                                           per student    per student\par
\noindent lab 3                                                                            (sheets.\par
\noindent google.com)\par
\noindent 6\par

\medskip\noindent\textbf{PROCEDURE}\par
\noindent Activity: Creating Different Types of Charts\par
\noindent 1. Open Google Sheets.\par
\noindent 2. Create a blank spreadsheet.\par
\noindent 3. Enter sample data.\par
\noindent 4, Select the dataset.\par
\noindent 5. Open Insert → Chart.\par
\noindent 6. Create:\par
\noindent\textbullet\; Column Chart\par
\noindent\textbullet\; Bar Chart\par
\noindent\textbullet\; Pie Chart\par
\noindent\textbullet\; Line Chart\par
\noindent\textbullet\; Scatter Chart\par
\noindent 7, Add titles and labels.\par
\noindent 8, Analyze the generated charts.\par
\noindent 9. Save the spreadsheet.\par

\medskip\noindent\textbf{RESULTS}\par
\noindent Thus different types of charts were created and analyzed successfully using a spreadsheet\par
\noindent application.\par
\clearpage
\medskip\noindent\textbf{PRACTICAL RELATED QUESTPRACTICAL RELATED QUESTPRACTICAL RELATED QUESTIONS 1,}\par
\noindent What is data visualization?\par
\noindent 2. What is a chart?\par
\noindent 3. What is the purpose of a Pie Chart?\par
\noindent 4. Which chart is used for trends?\par
\noindent 5. Which chart is suitable for comparison?\par
\noindent 6. What is a legend?\par
\noindent 7. What is the use of chart title?\par
\noindent 8, What is a Scatter Chart?\par
\noindent 9. Which application was used? 10, State one advantage of data visualization.\par

\medskip\noindent\textbf{SELF-LEARNING COMPONENTS}\par
\noindent 1. Create a small summary table showing the count of students who ‘Passed’ (Score > 40) versus\par
\noindent ‘Failed’ (Score < 40) in the Series Test of Foundational IT Skills. Insert a 3D Pie Chart to\par
\noindent represent these two groups as percentages of the whole class.\par
\noindent 2. Create a Clustered Bar Chart (Horizontal) that displays the students’ names on the Y-axis and their\par
\noindent marks on the X-axis. After creating the chart, apply a Chart Style from the Design tab to change\par
\noindent the background to a dark theme.\par
\noindent REFERENCES SI. No. Text Book/ Website ‘Author Publication 1 https://support.google.com/doc\par
\noindent Google Google Sheets Help\par
\noindent s/answer/63824 Centre\par

\noindent 2                           https://www.geeksforgee   GeeksforGeeks      — GeeksforGeeks data-vi\par
\noindent ks.org/                                      sualization/types-of-grap\par
\noindent hs-in-statistics/\par
\clearpage
\expheading{10}{\underline{\hspace{1.8cm}}}
\exptitle{AI SLIDE GENERATION}
\medskip\noindent\textbf{PRACTICAL SIGNIFICANCE}\par
\noindent To generate professional multi-slide presentations using Generative AI platforms by providing\par
\noindent structured text prompts.\par

\medskip\noindent\textbf{RELEVANT PROGRAM OUTCOMES(POS)}\par
\noindent PO1: Apply knowledge of mathematics, natural sciences, engineering fundamentals, and an\par
\noindent engineering specialization (as specified in DK1 to DK4) to a wide range of practical procedures and\par
\noindent practices\par
\noindent PO2: Identify and analyze well-defined engineering problems and reach substantiated conclusions\par
\noindent using codified methods of analysis specific to their field of activity (DK1 to DK4).\par
\noindent PO4: Conduct investigations of well-defined problems by locating and search relevant codes and\par
\noindent catalogues, and by conducting standard tests and measurements (DK8).\par
\noindent PO7: Understand and commit to professional ethics and the norms of technician practice, including\par
\noindent compliance with relevant laws. Demonstrate an understanding of the need for diversity and inclusion\par
\medskip\noindent\textbf{(DK9)}\par

\medskip\noindent\textbf{RELEVANT COURSE OUTCOMES(COS)}\par
\noindent CO3: Apply spreadsheet functions and Generative AI tools to construct interactive dashboards,\par
\noindent automated reports, and professional visual presentations for interpreting complex datasets.\par

\medskip\noindent\textbf{THEORY}\par
\noindent Generative AI\par
\noindent Generative AI is a type of Artificial Intelligence that creates new content such as text, images,\par
\noindent presentations, and reports based on user input.\par

\noindent Applications:\par
\noindent\textbullet\; Presentation generation\par
\noindent\textbullet\; Report generation\par
\noindent\textbullet\; Content creation\par
\noindent\textbullet\; Image generation\par
\noindent\textbullet\; Educational support Advantages:\par
\noindent\textbullet\; Saves time\par
\noindent\textbullet\; Improves productivity\par
\noindent\textbullet\; Generates organized content\par
\noindent\textbullet\; Supports creativity\par
\noindent AI slide generation refers to creating presentation slides automatically from user prompts. Basic\par
\noindent workflow:\par
\noindent 1. Enter prompt\par
\noindent 2. AIl generates slide structure\par
\noindent 3. AI creates content and design\par
\noindent 4. User edits output\par
\noindent 5. Presentation is exported\par
\noindent Prompt Engineering for Presentation Generation A structured prompt should contain:\par
\clearpage
\noindent\textbullet\; Topic\par
\noindent\textbullet\; Number of slides\par
\noindent\textbullet\; Target audience\par
\noindent\textbullet\; Main sections\par
\noindent\textbullet\; Style of presentation\par

\noindent Example Prompt:\par
\noindent "Create a 5-slide presentation on Cloud Computing for diploma students. Include introduction,\par
\noindent applications, advantages, challenges, and conclusion using professional style."\par

\medskip\noindent\textbf{RESOURCES REQUIRED}\par
\noindent Sl.No. [Name of resource               Computer System Desktop/Laptop    1 per student\par
\noindent Suggested specification Quantity 1     with Ubuntu\par

\noindent Linux 2                                Internet Connection Broadband /   Web Browser Mozilla Firefox /\par
\noindent Wi-Fi 1 per lab 3                 Google Installed\par

\noindent Chrome 4              Google Account         1 per student email   Gamma AI Accessed     Google Classroom\par
\noindent Valid institution or   ID 5                  via browser 1 per     Accessed via\par
\noindent personal                                     student               browser 1 per\par
\noindent (gamma.app) 6         student\par

\medskip\noindent\textbf{PROCEDURE}\par
\noindent Activity: Generating Slides using Gamma AI\par
\noindent 1. Open web browser in Ubuntu,\par
\noindent 2. Open Gamma AI.\par
\noindent 3. Sign in\par
\noindent 4. Create a new presentation.\par
\noindent 5. Enter a structured prompt.\par
\noindent 6. Generate outline.\par
\noindent 7. Generate slides.\par
\noindent 8. Review content.\par
\noindent 9. Modify slides if required.\par
\noindent 10, Export presentation.\par

\medskip\noindent\textbf{RESULTS}\par
\noindent A professional multi-slide presentation was successfully generated using a Generative AI platform by\par
\noindent providing structured text prompts. The generated presentation was reviewed,\par
\noindent AIN\par
\medskip\noindent\textbf{SCHEME}\par
\noindent refined, and customized as required. A report containing the prompt used, screenshots of the AI\par
\noindent generation process, screenshots of the final presentation was prepared and successfully uploaded\par
\noindent through Google Classroom.\par
\medskip\noindent\textbf{PRACTICAL RELATED QUESTPRACTICAL RELATED QUESTPRACTICAL RELATED QUESTIONS}\par
\noindent 1. What is Generative AI?\par
\noindent What is Gamma AI?\par
\noindent Ps\par
\noindent What is a prompt?\par
\clearpage
\noindent a\par
\noindent What is prompt engineering?\par
\noindent wo\par
\noindent Why should prompts be specific?\par
\noindent 6. Why should AI-generated content be reviewed?\par
\noindent 7. Name one AI presentation platform.\par
\noindent 8, What is the purpose of exporting presentations?\par
\noindent 9. What are the advantages of AI slide generation? 10, State one application of Generative AI.\par

\medskip\noindent\textbf{SELF-LEARNING COMPONENTS}\par
\noindent 1, Use Gamma.app to generate an 8-slide deck. Record the time taken (in seconds) for the AI to\par
\noindent generate the full presentation.\par
\noindent 2. Use Manus AI to research the current market price of three different laptop brands (Dell, HP, and\par
\noindent Lenovo) with 16GB RAM. Ask the agent to summarize these prices into a formatted table and then\par
\noindent create a 5-sentence summary recommending the best value for a student budget and make a\par
\noindent presentation.\par

\noindent SINo.                 Text Book/ Website    https:/help.gamma.   Gamma —\_         Gamma Official\par
\noindent ‘Author Publication   app/e                                 Help w/\par
\noindent 1                                                           Technologies Inc.\par

\noindent 2                           https:/vww.geeksforgee    GeeksforGeeks         GeeksforGeeks\par
\noindent ks.org/artificial-intelligen\par
\noindent ce/what-is-generative-ai/\par
\clearpage
\expheading{11}{\underline{\hspace{1.8cm}}}
\exptitle{GENERATIVE PRE-TRAINED}
\medskip\noindent\textbf{TRANSFORMER (GPT) GUIDED LEARNING}\par
\medskip\noindent\textbf{LEARNING}\par

\medskip\noindent\textbf{PRACTICAL SIGNIFICANCE}\par
\noindent To create and use a structured "Study Assistant" prompt in a Generative Pre-trained Transformer\par
\noindent (GPT) system for learning technical concepts effectively.\par

\medskip\noindent\textbf{RELEVANT PROGRAM OUTCOMES(POS)}\par
\noindent PO1: Apply knowledge of mathematics, natural sciences, engineering fundamentals, and an\par
\noindent engineering specialization (as specified in DK1 to DK4) to a wide range of practical procedures and\par
\noindent practices\par
\noindent PO2: Identify and analyze well-defined engineering problems and reach substantiated conclusions\par
\noindent using codified methods of analysis specific to their field of activity (DK1 to DK4).\par
\noindent PO3: Design solutions for well-defined technical problems and assist in the design of systems,\par
\noindent components, or processes to meet specified needs, with appropriate consideration ‘for public health\par
\noindent and safety, as well as cultural, societal, and environmental factors (DK5).\par
\noindent PO4: Conduct investigations of well-defined problems by locating and search relevant codes and\par
\noindent catalogues, and by conducting standard tests and measurements (DK8).\par
\noindent POT: Understand and commit to professional ethics and the norms of technician practice, including\par
\noindent compliance with relevant laws. Demonstrate an understanding of the need for diversity and inclusion\par
\medskip\noindent\textbf{(DK9).}\par

\medskip\noindent\textbf{RELEVANT COURSE OUTCOMES(COS)}\par
\noindent CO4: Apply generative AI, digital design platforms for guided learning and visual communication.\par

\medskip\noindent\textbf{THEORY}\par
\noindent Generative Pre-trained Transformer (GPT)\par
\noindent Generative Pre-trained Transformer (GPT) is an Artificial Intelligence model capable of\par
\noindent understanding and generating human-like text based on user input.\par

\noindent Applications:\par
\noindent © Educational assistance\par
\noindent © Question answering\par
\medskip\noindent\textbf{(TEACHERS AND ERY ~ DIPLOMA 2}\par
\noindent ME\par
\noindent Content generation\par
\noindent Programming support\par
\noindent Report writing\par
\noindent Language learning\par

\noindent Advantages:\par
\noindent Interactive learning\par
\clearpage
\noindent Instant explanations\par
\noindent Personalized assistance\par
\noindent Multiple examples and exercises\par

\medskip\noindent\textbf{Prompt Engineering}\par
\noindent Prompt engineering is the process of creating clear instructions to obtain better responses from AI\par
\noindent systems,\par

\noindent Characteristics of a good prompt:\par
\noindent Clearly specify the topic ‘Mention learning level\par
\noindent Specify expected output\par
\noindent Mention examples or questions required\par
\noindent Define response format\par
\noindent Study Assistant Prompt A Study Assistant prompt helps students learn topics systematically. Typical\par
\noindent structure:\par
\noindent Subjectiopic\par
\noindent Student level\par
\noindent Explanation requirement\par
\noindent Examples\par
\noindent Practice questions\par
\noindent Summary\par
\noindent aeAwNS\par

\noindent Sample Study Assistant Prompt\par
\noindent Prompt: "Act as a study assistant for diploma Computer Engineering students. Teach the topic Data\par
\noindent Structures ( Any branch specific topic) in simple language. Explain definitions, concepts, examples,\par
\noindent applications, important points, and provide five practice questions with answers."\par

\medskip\noindent\textbf{RESOURCES REQUIRED}\par
\noindent Sl.No.                      Name of resource            Computer System            1 per student\par
\noindent Suggested specification     Desktop/Laptop with\par
\noindent Quantity 1                  Ubuntu\par

\noindent Linux 2                                Internet Connection Broadband /     Web Browser Mozilla Firefox /\par
\noindent Wi-Fi 1 per lab 3                   Google Installed\par

\noindent Chrome 4           Google Account      Valid account    1 per student ID   GPT Platform     1 per student\par
\noindent / Email             for platform     5                  ChatGPT\par
\noindent login                               (chatgpt.com)/\par

\noindent Google Gemini / Any available\par

\noindent GPT tool 6                                              Google Classroom Accessed via browser 1 per\par
\noindent student\par

\medskip\noindent\textbf{PROCEDURE}\par
\noindent List the one-by-one procedure to be followed for completing the experiment.\par
\noindent Activity: Creating a Study Assistant Prompt\par
\noindent Open a web browser in Ubuntu.\par
\noindent Open GPT platform.\par
\clearpage
\noindent Sign in if required.\par
\noindent Create a structured prompt.\par
\noindent Enter prompt into the system.\par
\noindent Generate response.\par
\noindent Review the content.\par
\noindent Refine prompt if necessary.\par
\noindent Generate additional questions/examples. 10, Save observations.\par

\medskip\noindent\textbf{RESULTS}\par
\noindent A structured "Study Assistant” prompt was successfully created and used in a Generative Pre-trained\par
\noindent Transformer (GPT) system to learn technical concepts effectively. The GPT-generated explanations,\par
\noindent examples, summaries, and practice questions were analy:\par
\noindent and validated. A report containing the prompt used, generated responses, observations, and\par
\noindent necessary screenshots was prepared and successfully uploaded through Google Classroom.\par
\medskip\noindent\textbf{PRACTICAL RELATED QUESTPRACTICAL RELATED QUESTPRACTICAL RELATED QUESTIONS}\par
\noindent Expand GPT.\par
\noindent What is prompt engineering?\par
\noindent What is a Study Assistant prompt?\par
\noindent Why should prompts be structured?\par
\noindent Name one GPT platform.\par
\noindent What is the advantage of AI-assisted learning? Why should AI-generated content be reviewed? What\par
\noindent is prompt refinement?\par
\noindent State one application of GPT.\par
\noindent 10, What is the purpose of follow-up prompts?\par
\noindent SeNAwaAwNN\par
\noindent ‘SELF-LEARNING COMPONENTS List self-learning activities depending on the time allotted in the\par
\noindent syllabus.\par
\noindent 1. Give ChatGPT this prompt: "Act as a technical interviewer for a Junior Engineer role. Ask me 5\par
\noindent questions one by one about computers."\par
\noindent 2. What you mean by C-R-T-O framework in prompt engineering. Generate a proper prompt in the\par
\noindent C-R-T-O framework for use cases of your choice.\par
\noindent REFERENCES. SI.No. Text Book/ Website ‘Author Publication\par

\noindent 1                           https:/www.geeksforgeek   GeeksforGeeks      GeeksforGeeks\par
\noindent s.org/pro                                    mpt-engineering/\par

\medskip\noindent\textbf{(TEACHERS AND ERY ~ DIPLOMA 2}\par
\clearpage
\expheading{12}{\underline{\hspace{1.8cm}}}
\exptitle{GRAPHIC DESIGN}
\medskip\noindent\textbf{PRACTICAL SIGNIFICANCE}\par
\noindent To design a technical event poster using AI-assisted design features in Canva Free Version.\par
\noindent RELEVANT PROGRAM OUTCOMES(POS) PO1: Apply knowledge of mathematics, natural sciences,\par
\noindent engineering fundamentals, and an\par
\noindent engineering specialization (as specified in DK1 to DK4) to a wide range of practical procedures and\par
\noindent practices\par
\noindent P02: Identify and analyze well-defined engineering problems and reach substantiated conclusions\par
\noindent using codified methods of analysis specific to their field of activity (DK1 to DK4).\par
\noindent PO3: Design solutions for well-defined technical problems and assist in the design of systems,\par
\noindent components, or processes to meet specified needs, with appropriate consideration ‘for public health\par
\noindent and safety, as well as cultural, societal, and environmental factors (DK5).\par
\noindent PO4: Conduct investigations of well-defined problems by locating and search relevant codes and\par
\noindent catalogues, and by conducting standard tests and measurements (DK8).\par
\noindent PO7: Understand and commit to professional ethics and the norms of technician practice, including\par
\noindent compliance with relevant laws. Demonstrate an understanding of the need for diversity and inclusion\par
\medskip\noindent\textbf{(DK9)}\par

\medskip\noindent\textbf{RELEVANT COURSE OUTCOMES(COS)}\par
\noindent CO4: Apply generative AI, digital design platforms for guided learning and visual communication.\par

\medskip\noindent\textbf{THEORY}\par
\noindent Graphic design is the process of creating visual content to communicate information effectively using\par
\noindent text, images, colors, and layouts.\par
\noindent Applications.\par
\noindent Event posters\par
\noindent Advertisements\par
\noindent Social media posts\par
\noindent Technical brochures\par
\noindent © Presentation graphics\par
\noindent AI in Graphic Design Artificial Intelligence tools assist users in creating designs automatically by\par
\noindent suggesting:\par
\noindent Templates\par
\noindent Images\par
\noindent Layouts\par
\noindent Backgrounds\par
\noindent Design styles\par
\noindent Text content\par

\noindent Advantages:\par
\noindent Reduces design time\par
\noindent Provides professional templates\par
\clearpage
\noindent Improves creativity\par
\noindent Easy customization\par

\medskip\noindent\textbf{Canva}\par
\noindent Canva is an online graphic design platform used for creating visual content such as posters,\par
\noindent presentations, banners, and social media graphics.\par
\medskip\noindent\textbf{Features.}\par
\noindent Drag-and-drop interface\par
\noindent Templates\par
\noindent AI-assisted design tools\par
\noindent Icons and images\par
\noindent Text editing\par
\noindent Export options\par

\medskip\noindent\textbf{RESOURCES REQUIRED}\par
\noindent Sl.No.                                                   Name of resource Suggested specification ‘Quantity\par

\noindent 1                                      Computer System Desktop/Laptop      1 per student\par
\noindent with Ubuntu\par

\noindent Linux\par
\noindent Ps Internet Connection Broadband / Wi-Fi 1 per lab\par
\noindent 3 Web Browser Mozilla Firefox / Google Installed\par
\noindent Chrome\par

\noindent 4                                      Google Account ‘Valid institution   1 per student email ID\par
\noindent or personal\par

\noindent 5 Canva Free Version, accessed via 1 per student\par
\noindent browser (canva.com)\par

\noindent 6                                      College Logo File PNG format with   Google Classroom Accessed via\par
\noindent transparent 1 background            browser 1 per student\par
\noindent (provided by institution) 7\par

\medskip\noindent\textbf{PROCEDURE}\par
\noindent List the one-by-one procedure to be followed for completing the experiment.\par
\noindent Activity: Creating Technical Event Poster\par
\noindent 1. Open Canva Website Open the browser on your computer. Type www.canva.com in the address bar\par
\noindent and\par

\noindent press Enter:\par
\noindent 2. Sign in.\par
\noindent Click "Sign up" if new, or "Log in" if you already have an account, Click "Continue with Google" and\par
\noindent select your email ID. You will be taken to the Canva Home Screen.\par
\noindent 3. Select Poster design.\par
\noindent Click the search bar at the top, type "Poster" and press Enter. Select "Poster (Portrait)" from the\par
\noindent results. Click "Create a blank poster" to open the editor.\par
\noindent 4. Use AI prompt for poster generation.\par
\clearpage
\noindent In the left side panel, click "Magic Design” (starAvand icon). A text box will appear. Type a prompt\par
\noindent describing your event poster, for example:\par
\noindent "Technical event poster for a symposium on Artificial Intelligence, blue and white color scheme, with\par
\noindent space for event name, date, venue, and chief guest details."\par
\noindent Click "Generate". Wait a few seconds. Canva AI will display 4 to 6 ready-made poster designs. Click\par
\noindent the design you like best to open it in the editor.\par
\noindent Edit Event Details\par
\noindent Double-click on any text in the poster to edit it. Press Ctrl + A to select the existing text, delete it, and\par
\noindent type your own content. Replace all placeholder text with actual ‘event information:\par
\noindent Event name (e.g., TECHNOVA 2025)\par
\noindent Event type (e.g., National Level Technical Symposium)\par
\noindent Date, Time, and Venue\par
\noindent Chief Guest name and designation\par
\noindent Department and College name\par
\noindent Contact email or phone number.\par
\noindent 5. Edit Event Details.\par
\noindent Double-click on any text in the poster to edit it. Press Ctrl + A to select the existing text, delete it, and\par
\noindent type your own content. Replace all placeholder text with actual event information:\par
\noindent Event name (e.g., TECHNOVA 2025)\par
\noindent Event type (e.g., National Level Technical Symposium)\par
\noindent Date, Time, and Venue\par
\noindent Chief Guest name and designation\par
\noindent Department and College name\par
\noindent Contact email or phone number\par
\noindent 6. Format Text\par
\noindent Click on a text box and use the top toolbar to change the font style, font bold/talic, and text color. Use\par
\noindent large font size (60-80pt) for the event title and smaller size (I8-24pt) for body text. Ensure the text\par
\noindent color contrasts clearly with the background.\par
\noindent 7. Add Icons and College Logo\par
\noindent Click "Elements" in the left panel. Search for keywords like "technology" or "circuit" and click any\par
\noindent free icon to add it to the poster. To add your college logo, click "Uploads" — "Upload files”, select the\par
\noindent logo from your computer, and drag it to the top or bottom of the poster.\par
\noindent 8. Adjust Background\par
\noindent Click on the background area of the poster. Use the color box in the top toolbar to change the\par
\noindent background color. Choose a color that contrasts well with your text, such as dark blue with white\par
\noindent text.\par
\noindent 9. Preview the Poster\par
\noindent Click the eye icon (Preview) at the top-right corner. Check for spelling errors, overlapping elements,\par
\noindent and missing information. Press Escape to return to the editor and make corrections if needed.\par
\noindent 10. Download the Poster\par
\noindent Click "Share" at the top-right, then click "Download". Select PNG for digital submission or PDF for\par
\noindent printing. Click "Download" — the file will be saved in your Downloads folder.\par
\noindent II. Upload to Google Classroom\par
\clearpage
\noindent Open Google Classroom, go to the assigned practical, click "Add or Create" > "File", select the\par
\noindent downloaded poster file, and click "Turn In".\par

\medskip\noindent\textbf{RESULTS}\par
\noindent A technical event poster was successfully designed using AI-assisted design features available in\par
\noindent Canva Free Version. The generated design was customized by modifying text, layout, graphics, and\par
\noindent visual elements to meet the event requirements, A report containing the design prompt, the final\par
\noindent poster was prepared and successfully uploaded through Google Classroom.\par
\medskip\noindent\textbf{PRACTICAL RELATED QUESTPRACTICAL RELATED QUESTPRACTICAL RELATED QUESTIONS}\par
\noindent What is graphic design?\par
\noindent What is Canva?\par
\noindent What is AI-assisted design?\par
\noindent What is Magic Design?\par
\noindent Why are templates used?\par
\noindent What is the purpose of posters?\par
\noindent ‘Name two export formats in Canva.\par
\noindent What is the use of icons in posters?\par
\noindent Why should color combinations be selected carefully?\par
\noindent 10. State one advantage of AI in graphic design.\par

\medskip\noindent\textbf{SELF-LEARNING COMPONENTS}\par
\noindent 1, Use canva or any web based graphic design platform to design logo for your college art\par
\noindent festival\par
\noindent 2. Use Canva’s Magic Studio or any web based graphic design platform to place a photo of yourself in\par
\noindent front of a "High-Tech Server Room" background for a professional poster.\par

\noindent [REFERENCES           https:/\textbackslash\{\}www.canva.c   hups:/www.geeksfo    GeeksforGeeks   GeeksforGeeks\par
\noindent SINo. Text Book!      om/en\_in/help/use     rgeeks.org/graphic\par
\noindent Website ‘Author       Canva Canva\par
\noindent Publication 1         Official\par
\noindent -magic-design/ 2 \_\par

\noindent -design/\par
\medskip\noindent\textbf{STATE INS}\par
\medskip\noindent\textbf{[AND RESEARCH, KAIAMA}\par
\medskip\noindent\textbf{UTE OF TECHN}\par
\end{document}
